% !TeX root = ../main.tex
\ChapterOpening{cross-campus-training}{2}{跨校区训练车费报销}
\ChapterGuide
  {\GuideItem{\hyperref[invoice-title]{中山大学开票信息}}
   \GuideItem{\hyperref[evidence-chain]{各类凭证的用途}}
   \GuideItem{\hyperref[material-packaging]{票据整理与打包}}}
  {\GuideItem{\hyperref[cross-campus-training]{跨校区训练车费报销}}}
  {\GuideItem{\hyperref[training-allowance]{训练费发放}}}
\setlength{\parindent}{2\ccwd}

本章主要介绍跨校区训练车费的报销流程。对于不常驻广州校区南校园的队员，如果想利用周末或节假日到南校进行训练，
可以报销从自己校区到南校园的往返车费。

往返的交通方式尽量选择校车（例如珠海到广州的岐关车），其他出行方式不保证能报销，笔者目前没有实例可以验证。



\Subheading{1. 校车发票}

校车订单完成后，可以在历史订单中选择开具发票。同样地，发票抬头和纳税人识别号应填写为中山大学及其对应识别号。
发票开具后，务必核对发票信息是否正确。

发票的开具时间有限制，请在行程完成之后及时申请发票。

\Subheading{2. 汇总表及情况说明}

在纸上标注你的姓名学号，发票张数，用加号隔开每张发票的金额以及总金额，同时如果有单程则需要说明单程原因。

e.g.\hspace{0.4cm} 从珠海到广州后直接从广州一起去赛场，比赛结束后直接从赛场回珠海，广州到赛场以及赛场回珠海的费用在比赛报销中，不归入跨校区训练车费报销

下面是参考模板，你可以\href{https://github.com/PhSeCl/SYSU-Reimbursement-Handbook/releases/download/v2026.11.05/cross-campus-summary-form.pdf}{\textcolor{BookBlue}{点击这里}}直接下载可填写的表单，也可以自己完成。

\DocumentFigure[.7\linewidth][colframe=black!25,colback=white,drop fuzzy shadow=black!25]{cross-campus-summary.pdf}{2.1}{跨校区训练车费汇总表及情况说明示例}

将以上文件打包直接交给负责教练即可。

尽量在\strong{年底封账前提交}今年的跨校区训练车费报销，否则第二年年初再提交可能因逾期而无法报销。

\begin{SideNote}
这个过程极其缓慢，可以多向教练询问进度，不然可能会被一直拖着超出报销期限。（只能报销今年的费用）\\
笔者亲身经历，甚至重复提交了部分发票文件，但还是被拖没了几百块钱，由于体育部处理太慢，导致发票过期无法报销
想要解决这个问题，你可以尝试以下两点：
\begin{itemize}
  \item 材料交给教练后，注意自己的企业微信的财务信息系统是否有更新，是否可以看到这一笔报销的进度；如果没有，那说明大概率还没有制单
  ，可以催促要求尽快制单
  \item 看到报销单进度后，如果财务部对此单有问题希望补充材料，千万不要发给教练或其他方式转交，请直接在系统找到财务部对应负责人，
  直接发送邮件或者企业微信联系，这样效率快得多。（财务部的处理效率相当之高，反馈及时的话甚至可以在一天内解决）
\end{itemize}
\end{SideNote}

\DocumentFigure[.6\linewidth]{training-notification.png}{2.2}{企业微信财务信息系统的接单通知}

\DocumentFigure{training-email-inquiry.png}{2.3}{就票据查验及报销期限问题向审核人员说明情况}

\DocumentFigure{training-email-resolution.png}{2.4}{收到审核反馈后，确认核减金额及后续处理方式}

\DocumentFigure{training-system-result.png}{2.5}{系统中的金额调整、处理进度与流转记录}
